% Modelo criado com base no Modelo Canônico de Artigo Científico com abnTeX2
% de abnTeX2 group
%
% Foram feitas mudanças significativas devido a necessidade de formatação
% deste tipo de documento
\documentclass[11pt]{article}
\usepackage[portuges]{babel}    % traduz elementos padrão como Abstract para
                                % Resumo
\usepackage{fontspec}           % permite alterar a font do texto, mas precisa 
                                % mudar o compilador do Overleaf para o XeLaTeX
                                % (menu >> compiler)
\usepackage{graphicx}           % permite inserir figura no texto
\usepackage{hyperref}           % permite inserir links no texto
\usepackage{sectsty}            % para configurar o tamanho da fonte das seções
\usepackage{color}				% Controle das cores
                                % dos acentos)
\usepackage{microtype} 			% para melhorias de justificação
\usepackage{geometry}           % permite manipular o tamanho da página e margens
%-------------------------------------------------------------------------------
% Usado para gerar as tabelas necessárias para o cronograma, bem como as
% configurações necessárias para os alinhamentos corretos nas células
%-------------------------------------------------------------------------------
\usepackage{multirow}           % permite criar tabelas com celulas agrupadas
\usepackage{tabularx}           % permite quebra automática de linhas em tabela
\renewcommand\tabularxcolumn[1]{m{#1}}% for vertical centering text in X column
%-------------------------------------------------------------------------------
% Pacotes de citações
%-------------------------------------------------------------------------------
\usepackage[brazilian,hyperpageref]{backref}	 % Paginas com as citações na bibl
\usepackage[alf]{abntex2cite}	% Citações padrão ABNT


%-------------------------------------------------------------------------------
% informações do PDF
%-------------------------------------------------------------------------------
% alterando o aspecto da cor azul
\definecolor{blue}{RGB}{41,5,195}
\makeatletter
\hypersetup{
		pdftitle={\@title}, 
% 		pdfauthor={\@author}, %autor nao pode ser identificado
    	pdfsubject={Modelo de Projet PIBIT},
	    pdfcreator={Joaquim Pessôa Filho},
		pdfkeywords={abnt}{latex}{mackenzie}{projeto}{iniciação científica}, 
		colorlinks=true,       		% false: boxed links; true: colored links
    	linkcolor=blue,          	% color of internal links
    	citecolor=blue,        		% color of links to bibliography
    	filecolor=magenta,      		% color of file links
		urlcolor=blue,
		bookmarksdepth=4
}
\makeatother
%-------------------------------------------------------------------------------
% Configuração da Font e espaçamentos seguindo a norma estabelecida para
% submissão do projeto
%-------------------------------------------------------------------------------
\sectionfont{\fontsize{11}{15}\selectfont}  % Fonte Size dos títulos das seções 
                                            % em 11pt
\setmainfont{Arial}                         % todo o texto em Arial
\setlength{\parindent}{0pt}                 % tabulação no início dos parágrafos
\setlength{\parskip}{6pt}                   % espaço entre parágrafos
\renewcommand{\baselinestretch}{1.5}        % espaço entre linhas
\geometry{
    a4paper,                                % tamanho A4
    total={160mm,257mm},                    % tamanho do texto na página
    left=30mm,                              % margem esquerda
    top=20mm,                               % margem do topo
}
%-------------------------------------------------------------------------------
% Configurações do pacote backref
%-------------------------------------------------------------------------------
% Usado sem a opção hyperpageref de backref
\renewcommand{\backrefpagesname}{Citado na(s) página(s):~}
% Texto padrão antes do número das páginas
\renewcommand{\backref}{}
% Define os textos da citação
\renewcommand*{\backrefalt}[4]{
	\ifcase #1 %
		Nenhuma citação no texto.%
	\or
		Citado na página #2.%
	\else
		Citado #1 vezes nas páginas #2.%
	\fi}%
%-------------------------------------------------------------------------------
% Muda a orientação do título para esquerda e deixa todas as letras em maiúscula
%-------------------------------------------------------------------------------
\makeatletter
\renewcommand{\maketitle}{\bgroup\setlength{\parindent}{11pt}
\begin{flushleft}
  \MakeUppercase{\textbf{\@title}}
\end{flushleft}\egroup
}
\makeatother
%-------------------------------------------------------------------------------
% Muda a formatação do abstract
%-------------------------------------------------------------------------------
\renewenvironment{abstract}
 {\par\noindent\textbf{\abstractname}\\ \ignorespaces}
 {\par\medskip}



\title{Insira o Título do seu Projeto}
% \author{**** NÃO INCLUIR NOMES DOS AUTORES NEM A UNIDADE ACADÊMICA ***}


\begin{document}

\maketitle

\begin{abstract}
(Resumo do projeto. Motivação, justificativa, relevância, objetivo principal, resultados
esperados.)
\end{abstract}

\section{Objetivos}
(Enumere e detalhe o objetivo geral e objetivos específicos do projeto. Projetos de iniciação
tecnológica podem estar relacionados a atividades tais como: desenvolvimento de protótipos,
produtos, processos, métodos e técnicas; levantamentos de informação tecnológica;
mapeamento de público-alvo e mercado; avaliação de soluções em laboratório ou sua
observação de uso por usuários finais; desenvolvimento de atividades de capacitação e/ou
incentivo à inovação, empreendedorismo e transferência de tecnologia ao mercado ou à
sociedade, entre outros objetivos.


\section{Cenário Problema}
(Descreva o problema que o projeto busca resolver ou reduzir. Caracterize o contexto em que
o projeto está inserido em termos de domínio de aplicação, público-alvo, barreiras,
dificuldades, impactos causados pelo problema ou outras informações que julgar relevantes
para o entendimento do problema.)


\section{Estado da Arte e da Prática}
(Apresente os principais conceitos utilizados no projeto, com base em referências científicas
e/ou tecnológicas. Relate os avanços mais relevantes do estado da arte (pesquisa) e da prática
(mercado/sociedade) de soluções para o problema abordado no projeto. Cuide para que não
seja uma simples descrição revisional de textos ou fragmentos de texto que não dialogam entre
si. Ao contrário, esta seção deve apresentar uma estrutura coesa e coerente com o problema
e a solução a ser apresentada no projeto.)


\section{Desafio Tecno-científico}
(Apresente os desafios tecno-científicos a serem abordados pelo desenvolvimento do projeto,
considerando o estado da arte a respeito do problema tecno-científico a ser resolvido e o
estado da prática das soluções existentes atualmente em uso no mercado, organizações ou
sociedade.)


\section{Proposta de solução}
(Descreva as características principais da solução a ser desenvolvida no projeto e argumente
sobre seu potencial para solução do problema mencionado anteriormente.)


\section{Procedimentos Metodológicos}
(Descreva os passos para o desenvolvimento do projeto. Os procedimentos metodológicos
devem responder a pergunta: “Como este projeto será desenvolvido?”. Detalhe o suficiente
para que o caminho a ser percorrido no projeto seja compreendido.)


\section{Cronograma de Atividades}
(Distribua os passos para o desenvolvimento do projeto ao longo do tempo, usando a Tabela \ref{tab:cronograma}. Os projetos de PIBIT são realizados, normalmente, no prazo de 12 meses. Portanto, o
cronograma precisa ser coerente com esse prazo. Algumas vezes os objetivos precisarão ser
readequados para cumprir o prazo proposto. É muito importante que o aluno e o professor-
orientador verifiquem se a proposta de projeto será suficientemente atendida no prazo dos 12
meses.)


\begin{table}[h]
\centering
\begin{tabularx}{\textwidth}{|X|cccccccccccc|}
\hline
\multicolumn{1}{|c|}{\multirow{2}{*}{Atividades}} &
    \multicolumn{12}{c|}{Meses} \\ \cline{2-13} 
 &
    \multicolumn{1}{c|}{1} &
    \multicolumn{1}{c|}{2} &
    \multicolumn{1}{c|}{3} &
    \multicolumn{1}{c|}{4} &
    \multicolumn{1}{c|}{5} &
    \multicolumn{1}{c|}{6} &
    \multicolumn{1}{c|}{7} &
    \multicolumn{1}{c|}{8} &
    \multicolumn{1}{c|}{9} &
    \multicolumn{1}{c|}{10} &
    \multicolumn{1}{c|}{11} &
    12 \\ \hline
Atividade 1 &
    \multicolumn{1}{c|}{X} & % mês 1
    \multicolumn{1}{c|}{ } & % mês 2
    \multicolumn{1}{c|}{ } & % mês 3
    \multicolumn{1}{c|}{ } & % mês 4
    \multicolumn{1}{c|}{ } & % mês 5
    \multicolumn{1}{c|}{ } & % mês 6
    \multicolumn{1}{c|}{ } & % mês 7
    \multicolumn{1}{c|}{ } & % mês 8
    \multicolumn{1}{c|}{ } & % mês 9
    \multicolumn{1}{c|}{ } & % mês 10
    \multicolumn{1}{c|}{ } & % mês 11
    X \\ \hline % mês 12
Atividade 2 &
    \multicolumn{1}{c|}{X} & % mês 1
    \multicolumn{1}{c|}{X} & % mês 2
    \multicolumn{1}{c|}{X} & % mês 3
    \multicolumn{1}{c|}{ } & % mês 4
    \multicolumn{1}{c|}{ } & % mês 5
    \multicolumn{1}{c|}{ } & % mês 6
    \multicolumn{1}{c|}{ } & % mês 7
    \multicolumn{1}{c|}{ } & % mês 8
    \multicolumn{1}{c|}{ } & % mês 9
    \multicolumn{1}{c|}{ } & % mês 10
    \multicolumn{1}{c|}{ } & % mês 11
    X \\ \hline % mês 12
Atividade 3 &
    \multicolumn{1}{c|}{ } & % mês 1
    \multicolumn{1}{c|}{ } & % mês 2
    \multicolumn{1}{c|}{X} & % mês 3
    \multicolumn{1}{c|}{X} & % mês 4
    \multicolumn{1}{c|}{X} & % mês 5
    \multicolumn{1}{c|}{X} & % mês 6
    \multicolumn{1}{c|}{ } & % mês 7
    \multicolumn{1}{c|}{ } & % mês 8
    \multicolumn{1}{c|}{ } & % mês 9
    \multicolumn{1}{c|}{ } & % mês 10
    \multicolumn{1}{c|}{ } & % mês 11
    X \\ \hline % mês 12
 Caso tenha uma atividade com um título grande o tabularx irá quebrar automaticamente o texto &
    \multicolumn{1}{c|}{ } & % mês 1
    \multicolumn{1}{c|}{ } & % mês 2
    \multicolumn{1}{c|}{ } & % mês 3
    \multicolumn{1}{c|}{ } & % mês 4
    \multicolumn{1}{c|}{ } & % mês 5
    \multicolumn{1}{c|}{ } & % mês 6
    \multicolumn{1}{c|}{ } & % mês 7
    \multicolumn{1}{c|}{ } & % mês 8
    \multicolumn{1}{c|}{ } & % mês 9
    \multicolumn{1}{c|}{ } & % mês 10
    \multicolumn{1}{c|}{ } & % mês 11
    X \\ \hline % mês 12
 ... &
    \multicolumn{1}{c|}{ } & % mês 1
    \multicolumn{1}{c|}{ } & % mês 2
    \multicolumn{1}{c|}{ } & % mês 3
    \multicolumn{1}{c|}{ } & % mês 4
    \multicolumn{1}{c|}{ } & % mês 5
    \multicolumn{1}{c|}{ } & % mês 6
    \multicolumn{1}{c|}{ } & % mês 7
    \multicolumn{1}{c|}{ } & % mês 8
    \multicolumn{1}{c|}{ } & % mês 9
    \multicolumn{1}{c|}{ } & % mês 10
    \multicolumn{1}{c|}{ } & % mês 11
    X \\ \hline % mês 12
 ... &
    \multicolumn{1}{c|}{ } & % mês 1
    \multicolumn{1}{c|}{ } & % mês 2
    \multicolumn{1}{c|}{ } & % mês 3
    \multicolumn{1}{c|}{ } & % mês 4
    \multicolumn{1}{c|}{ } & % mês 5
    \multicolumn{1}{c|}{ } & % mês 6
    \multicolumn{1}{c|}{ } & % mês 7
    \multicolumn{1}{c|}{ } & % mês 8
    \multicolumn{1}{c|}{ } & % mês 9
    \multicolumn{1}{c|}{ } & % mês 10
    \multicolumn{1}{c|}{ } & % mês 11
    X \\ \hline % mês 12
\end{tabularx}
\caption{Cronograma de Atividades}
\label{tab:cronograma}
\end{table}


\section{Resultados esperados}
(Descreva o potencial do projeto para criação e disponibilização de novos produtos/processos,
propriedade intelectual, divulgação tecnológica, capacitação profissional, incentivo à inovação
e empreendedorismo, entre outros resultados possíveis.)


\section{Cenário de Aplicação e Impactos}
(Descreva como seria a aplicação da solução inovadora proposta. Por exemplo, como seria
usada por empresas, pessoas/consumidores, setor público ou terceiro setor, Ressalte os
impactos - sociais, ambientais, econômicos - associados a essa aplicação.)


\section{Nível de Maturidade Tecnológica}
(Descreva o nível de maturidade tecnológica da solução a ser desenvolvida usando a escala
TRL1 - Technology Readiness Level - ver Figura \ref{fig:trl}) \cite{johnc.mankins1995}:

\begin{figure}[h]
    \centering
    \includegraphics[scale=0.5]{imagens/trl.png}
    \caption{TRL - Technology Readiness Level.}
    \label{fig:trl}
\end{figure}


\section{Informações Adicionais}
(Acrescente quaisquer informações que gostaria de acrescentar sobre o projeto. Caso o projeto
envolva a parceria com empresas, apresente carta de interesse de empresa no projeto.)

\textbf{Referências}

(Enumere as referências importantes relacionadas ao projeto no formato ABNT. Lembre-se que referências dos próprios autores devem ser omitidas devido à avaliação cega requerida pelo edital. Inclua referências a fontes de informação tecnológica/propriedade intelectual relacionadas ao projeto: patentes, registro de software etc, se houver. No link a seguir é
possível encontrar as regras para cada tipo de referência, diretamente, \href{https://www.normasabnt.org/normas-abnt-2022/}{no site da ABNT (normas atualizadas 2022)} \cite{abnt-norma-2022})

\textbf{*** TAMANHO MÍNIMO DE 6 PGS E MÁXIMO DE 10 PGS ****}

% ----------------------------------------------------------
% Referências bibliográficas
% ----------------------------------------------------------
\bibliography{referencias}

\end{document}
